\begin{answer}

\subsubsection*{i)} 
A trivial classifier can either guess always positive or negative. Regardless of the dataset, there exists a classifier with either $A=p$ or $A=1-p$. 

For both cases, from a classifier with $p \le 1-p$ or $p > 1-p$, we can say that our accuracy could be at least $1-p$.


\subsubsection*{For $p > 1-p$:}
A trivial classifier with either have $A=p$ or $A=1-p$.
Due to $p>1-p$, we can say that $A=p$ or $A=1-p$ which satisfies $A \ge 1-p$

\subsubsection*{For $1-p > p$:}
A trivial classifier once again will either have $A=p$ or $A=1-p$.

Hence due to $1-p \ge p$ we have:

    i) $A=p \rightarrow A < 1-p \rightarrow$ A always positive classifier is not fulfilling \\
    ii) $A=1-p \rightarrow A > 1-p \rightarrow$ Satisfies $A \ge 1-p$.

$\Rightarrow$ For both cases, where $p>1-p$ or $p<=1-p$ there exists one classifier that satisfies $A \ge 1-p$.


\subsubsection*{ii)} The amount of all examples is the sum of positive as well as negative examples
\#examples = TP+FN+TN+FP

The amount of only positive examples is \#positive examples = TP+FN. \\
p is defined as the fraction of positive examples. Hence, it follows:
\[ p = \frac{\#\text{positive examples}}{\#\text{examples}} = \frac{TP+FN}{TP+FN+TN+FP} \]


From before, we know: \\
\#examples = TP+FN+TN+FP
And correct predictions examples are: \\
\#correct predictions = TP+TN

So, it follows: \\
\[ A = \frac{TP+TN}{TP+FN+TN+FP} \]
\[ A = \frac{TP}{TP+FN+TN+FP} + = \frac{TN}{TP+FN+TN+FP} \]

Let $N = TP+FN+TN+FP$: \\
\[ A = \frac{TP}{N} + \frac{TN}{N} \]
\begin{align*}
    A &= \underbrace{\frac{TP}{TP+FN}}_{A_1} \cdot \underbrace{\frac{TP+FN}{N}}_{p} + \underbrace{\frac{TN}{TN+FP}}_{A_2} \cdot \underbrace{\frac{TN+FP}{N}}_{(1-p)} \\
    &= A_1 p + A_2 (1-p)
\end{align*}

iii) The trivial classifier that always predicts negative, will have $A = 1-p$.

It's accuracy on positive examples $A_1$ is zero, as we always predict negative.

% A second side note is written below the first one.
The same way our accuracy $A_0$ is always 1.

This gives us:
$\bar{A} = \frac{1}{2}(1+0) = 0.5$


\end{answer}